\bookchapter{附錄H：版本路線圖}

\chapterintro{此書猶續有更新。下陳未來版本規劃之方向，亦歡迎爾於反饋中告吾儕所最欲觀之內容。}

\textbf{v1.8 已落地：新增第七、八、九篇「AI編程實戰篇」「安卓開發實戰篇」「iOS開發實戰篇」（第37-69章，共33章）}\par

自 AI 編程之範式革命、五大思維原則、上下文與需求工程、2026工具與模型全景、API實戰、對話工程、三個完整項目、質量保障與進階之路——攜爾自「會用AI」升級至「指揮AI寫代碼」；至安卓與iOS兩大平臺完整開發之鏈路：環境搭建（Android Studio / Xcode）、語言速成（Kotlin / Swift）、界面框架（Jetpack Compose / SwiftUI）、數據與網絡、實戰項目（記賬本 + 天氣查詢）、測試與上架（Google Play / App Store）及速查手冊——攜爾自數碼用戶終成 App 開發者。

\begin{tcbinfo}{下一版（v1.9）計劃方向}
\begin{itemize}
  \item 手機與可穿戴設備：Apple Watch、華為手錶、智能戒指（Oura/Galaxy Ring）深度協同指南
  \item 旅行與手機：eSIM實操、國際流量方案、離線地圖、海外支付、翻譯App對比
  \item 手機與健康管理：運動健康、睡眠監測、心率血氧、女性健康、心理健康App
  \item AI短視頻創作：剪映AI、可靈AI、數字人、AI配音之實操教程
  \item 豎屏短劇入門：劇本結構、拍攝技巧、平臺政策
\end{itemize}

\end{tcbinfo}

\begin{tcbwarn}{關於時效性之說明}
手機技術迭代極快，本書內容據2026年中技術與市場之況而撰。AI 之能力、芯片型號、系統版本、價格區間，爾閱讀之時或已變矣。涉具體操作步驟，請以爾設備實際系統版本為準；涉購買決策，請查閱最新評測與用戶反饋。

\end{tcbwarn}
